\section*{Chapitre \ref{ch:ll:cpp-for-latency-ii-compile-time} --- C++ pour la latence II~: temps de compilation}

\begin{solution}{exo:ll:cpp-for-latency-ii-compile-time:1}
Sur la mesure retenue, $1.94/0.52 \approx 3.7$ fois pour un seul type de message et $8.69/5.22 \approx 1.7$ fois sur le
mélange.
\end{solution}

\begin{solution}{exo:ll:cpp-for-latency-ii-compile-time:2}
$(7.26 - 1.14)\ \text{ns} \times 5 \times 10^6 \approx 31$ millisecondes de traitement par seconde~: 3~\% d'un cœur, et six
nanosecondes de latence ajoutées à chaque événement.
\end{solution}

\begin{solution}{exo:ll:cpp-for-latency-ii-compile-time:3}
$250 \times 10^2 = 25\,000$~; $2\,500 / 10^2 = 25$. \texttt{rescale(1, 2, 20)}~: l'échelle cible est hors plage, donc la
branche improbable appelle la fonction de rapport froide et renvoie 0.
\end{solution}

\begin{solution}{exo:ll:cpp-for-latency-ii-compile-time:4}
Le \texttt{switch} est compilé en table de sauts (un saut indirect) ou en chaîne de branchements conditionnels~; les deux dépendent du
type de message, et sur un mélange aléatoire le prédicteur se trompe environ trois fois sur quatre.
\end{solution}

\begin{solution}{exo:ll:cpp-for-latency-ii-compile-time:5}
L'étape~1 (ajout de la quantité) est le chargement de \texttt{qty} et le \texttt{add}~; l'étape~2 est le \texttt{xor} avec le prix et le
stockage de \texttt{acc}~; l'étape~3 est le \texttt{cmp} avec la limite et le \texttt{jge}~; l'étape~4 est le \texttt{add} à la somme en
mémoire. Les objets des étapes~1 et~2 ont disparu~: ils n'ont pas de membres de données, seulement du code, qui a été inliné~; la
limite de l'étape~3 et la somme de l'étape~4 subsistent sous forme de chargements et d'écritures aux décalages de l'objet pipeline.
\end{solution}

\begin{solution}{exo:ll:cpp-for-latency-ii-compile-time:6}
Faire du moteur un template sur le type de stratégie, et l'instancier pour chacune des trois~; au démarrage, lire la
configuration et appeler le \texttt{run} de l'instanciation correspondante une seule fois. La répartition unique a lieu hors de la
boucle critique.
\end{solution}

\begin{solution}{exo:ll:cpp-for-latency-ii-compile-time:7}
Le listing gagne un \texttt{call} vers l'étape, les champs de l'événement étant écrits en mémoire avant lui et rechargés après, puisque
le compilateur ne voit plus ce que la fonction appelée modifie~; le temps par événement remonte vers celui du pipeline dynamique, sans le
chargement indirect d'une table.
\end{solution}

\begin{solution}{exo:ll:cpp-for-latency-ii-compile-time:8}
Le \omterm{def:ll:cpu-microarchitecture:ubench}{microbenchmark} mesure un appel parfaitement prédit dans une boucle serrée. Il ne voit pas les erreurs de prédiction lorsque les cibles varient, et
surtout l'optimisation perdue~: un appel virtuel ne peut pas être inliné, donc les valeurs sont sauvegardées autour de lui et rien ne se simplifie
à travers. Le pipeline à quatre étapes a perdu six nanosecondes, et non pas quatre appels valant chacun 1.9.
\end{solution}

\begin{solution}{pb:ll:cpp-for-latency-ii-compile-time:1}
\begin{enumerate}
\item Un seul type~: virtuel 1.9, variant 1.8, CRTP \qty{0.5}{\nano\second}~; mélange~: 8.7, 6.4 et \qty{5.2}{\nano\second} (mesure
retenue).
\item \omterm{def:ll:cpu-microarchitecture:branch}{Erreur de prédiction de branchement} sur le type de message, qui choisit la cible.
\item Environ 1.1 et \qty{7.3}{\nano\second} par événement.
\item Ses étapes s'exécutent dans le même ordre pour chaque événement~; le seul branchement dépendant des données est la porte, identique dans les deux cas.
\item Il charge le pointeur vers l'étape suivante, charge la table de l'objet, appelle à travers elle, teste le résultat et boucle~; les registres
sont sauvegardés autour de chaque appel.
\item Treize instructions.
\item Replier ensemble l'arithmétique des étapes, garder l'événement dans les registres, supprimer les objets d'étape vides, et un
unique branchement conditionnel pour la porte.
\item Il permet au compilateur de dévirtualiser les appels par références du type concret~; il ne fait rien pour les appels passant par
le pointeur vers la classe de base qu'utilise une architecture à modules d'extension.
\item Faire de chaque module d'extension un paramètre de template et instancier le moteur pour chaque combinaison prise en charge~; choisir
l'instanciation une seule fois au démarrage.
\item Un nouveau module d'extension ne peut pas être chargé dans un binaire en cours d'exécution~; en ajouter un exige une reconstruction et une mise en production.
\item $3 \times 2 = 6$.
\item Chaque variante est construite et exécute la suite complète de tests et de rejeu~; une matrice dans l'intégration continue.
\item \textbf{Résultat à retenir.} La répartition virtuelle coûte environ 3.7 fois le CRTP lorsque le type de message est prévisible et environ 1.7
fois sur un mélange aléatoire~; le pipeline à quatre étapes coûte environ \qty{1.1}{\nano\second} par événement composé statiquement et
\qty{7.3}{\nano\second} via des appels virtuels.
\item Le \omterm{def:ll:cpu-microarchitecture:ubench}{microbenchmark} mesure l'appel seul, prédit, et non ce que cet appel empêche~: l'\omterm{def:ll:cpp-for-latency-ii-compile-time:inline}{inlining} et les optimisations qu'il rend impossibles.
\item Hors du chemin critique (hot path)~: configuration, rapports, outils, et là où l'ensemble des types est ouvert (modules d'extension chargés sans
reconstruction dans un moteur de recherche quantitative).
\item Une exigence nommée, vérifiée au point d'utilisation~: l'erreur pour une étape sans \texttt{bool on(E\&)} tient sur une seule ligne.
\item Tables et constantes (échelles, longueurs, tables de correspondance) construites sans code de démarrage et vérifiées par
\texttt{static\_assert}.
\item Garder la fonction critique petite et contiguë dans le cache d'instructions, le cas courant étant placé à la suite.
\item Parce que \texttt{always\_inline} est une demande adressée à une heuristique susceptible de changer d'une version de compilateur à l'autre~; le code
assembleur montre ce qui s'est passé.
\item Tout ce qui est fixé pour toute la durée du processus~: types, stratégies, tables.
\end{enumerate}
\end{solution}

\begin{solution}{iq:ll:cpp-for-latency-ii-compile-time:1}
L'objet contient un pointeur vers la table de pointeurs de fonction de sa classe~; l'appel charge ce pointeur, charge l'entrée et appelle
par son intermédiaire. Coût~: deux chargements dépendants, un branchement indirect (peu coûteux lorsqu'il est prédit), et la perte de l'\omterm{def:ll:cpp-for-latency-ii-compile-time:inline}{inlining}.
\iqlookfor{le mécanisme de la vtable et le coût indirect.}
\end{solution}

\begin{solution}{iq:ll:cpp-for-latency-ii-compile-time:2}
Une classe de base template paramétrée par sa classe dérivée appelle les fonctions de la classe dérivée par
\texttt{static\_cast<Derived\&>(*this)}~; la répartition est résolue à la compilation et peut être inlinée. On l'utilise pour une famille fixe de
comportements sur un chemin critique (gestionnaires, politiques).
\iqlookfor{répartition à la compilation et son compromis~: aucune substitution à l'exécution.}
\end{solution}

\begin{solution}{iq:ll:cpp-for-latency-ii-compile-time:3}
Une fonction \texttt{constexpr} peut être évaluée à la compilation lorsque ses arguments sont des constantes, et s'exécute sinon à
l'exécution~; une fonction \texttt{consteval} doit être évaluée à la compilation, si bien qu'un appel qui ne peut pas l'être est une erreur.
\iqlookfor{pouvoir contre devoir.}
\end{solution}

\begin{solution}{iq:ll:cpp-for-latency-ii-compile-time:4}
Il expose le corps de la fonction appelée aux optimisations de l'appelant~: les constantes se propagent, les valeurs restent dans les registres, le travail mort
disparaît. Il peut aussi ralentir le code en le faisant grossir au-delà du cache d'instructions ou du tampon de boucle, ou en inlinant du code rare
dans une boucle critique.
\iqlookfor{optimisation à travers la frontière, et taille du code.}
\end{solution}

\begin{solution}{iq:ll:cpp-for-latency-ii-compile-time:5}
Un variant~: l'ensemble des types de messages est fermé, l'objet ne demande aucune allocation sur le tas, la répartition est un saut sur l'indice, et
les gestionnaires peuvent être inlinés. Une hiérarchie seulement si l'ensemble doit être ouvert.
\iqlookfor{ensembles de types fermés contre ouverts, et allocation.}
\end{solution}

\begin{solution}{iq:ll:cpp-for-latency-ii-compile-time:6}
Le moteur est un template sur sa stratégie (et d'autres politiques)~; chaque stratégie prise en charge est une instanciation compilée dans
le binaire~; la configuration sélectionne une instanciation au démarrage, par une répartition unique hors de la boucle~; la boucle critique
est entièrement statique. Testée sous forme de matrice.
\iqlookfor{une seule répartition au démarrage, statique ensuite.}
\end{solution}
